%% 第 9 章：幂零算子与零空间链
%% 本文件由 tools/md2tex.py 自动生成，请勿直接编辑。

\chapter{幂零算子与零空间链}


\begin{jdefn}{幂零算子}{r:10:1}
若存在正整数 $q$，使：

\[
N^q=0,
\]

则称 $N$ 是幂零算子。

幂零算子反复作用足够多次后，会将整个空间中的所有向量映为零。
\end{jdefn}

\begin{jprop}{零空间链递增}{r:10:2}
对任意算子 $N$：

\[
\ker N
\subseteq
\ker N^2
\subseteq
\ker N^3
\subseteq\cdots.
\]
\end{jprop}

\begin{jproof}{证明思路}
若：

\[
N^k(v)=0,
\]

则再施加一次 $N$ 后仍为零：

\[
N^{k+1}(v)=0.
\]

因此：

\[
v\in\ker N^{k+1}.
\]
\end{jproof}

\begin{jprop}{零空间链稳定}{r:10:3}
若：

\[
\ker N^r=\ker N^{r+1},
\]

则：

\[
\ker N^r=\ker N^{r+1}=\ker N^{r+2}=\cdots.
\]
\end{jprop}

\begin{jproof}{证明思路}
假设某个向量经 $N^{r+2}$ 后为零。

那么它经一次 $N$ 后属于 $\ker N^{r+1}$。由于：

\[
\ker N^{r+1}=\ker N^r,
\]

它经一次 $N$ 后实际上已属于 $\ker N^r$，因此原向量属于 $\ker N^{r+1}$。

所以再高一层不会产生新方向。之后可重复同样推理。
\end{jproof}

\begin{jprop}{有限维空间中的稳定性}{r:10:4}
若：

\[
\dim V=n,
\]

则零空间链最多在第 $n$ 步之前稳定。

因为每次严格增长都会使维数至少增加 $1$，而零空间维数不可能超过 $n$。
\end{jprop}

\begin{jprop}{幂零算子的幂}{r:10:5}
若 $N$ 是 $n$ 维空间上的幂零算子，则：

\[
N^n=0.
\]
\end{jprop}

\begin{jproof}{证明思路}
幂零意味着某一层零空间最终等于整个 $V$。

零空间链至多经过 $n$ 次严格增长便稳定，因此在第 $n$ 步时已经有：

\[
\ker N^n=V.
\]

这等价于：

\[
N^n=0.
\]
\end{jproof}
